\chapter{Réalisation et guide d'utilisation}
\label{ch:realisation}

\section{Stack technique déployée}

\begin{table}[H]
\centering
\caption{Technologies de production (\url{https://lexhub.space})}
\label{tab:stack}
\begin{tabular}{@{}ll@{}}
\toprule
\textbf{Couche} & \textbf{Choix réel} \\
\midrule
Interface & Next.js, TypeScript, Tailwind ; FR / AR \\
API & FastAPI, SQLAlchemy, Alembic \\
Données & PostgreSQL~16, MinIO (objets), Redis~7 \\
Tâches & Celery (OCR, e-mails, rappels) \\
Dictée & Web Speech API (Chrome / Edge) ; MediaRecorder $\rightarrow$ Whisper OpenAI (Firefox + Terminé) ; import audio \\
IA texte & Modèle de langage cloud (composition, révision pédagogue, synthèse), consignes anti-invention ; noms masqués avant l'IA \\
OCR & OCR cloud sur les pièces \\
Export & Word (.docx) officiel + PDF (RTL arabe, armoiries, en-tête cabinet) \\
Infra & Docker Compose, Caddy (TLS), VPS, images GHCR \\
Exploitation BDD & Adminer conteneurisé via tunnel SSH local ; pas d'exposition web publique \\
\bottomrule
\end{tabular}
\end{table}

La dictée combine deux voies : (1)~\textit{Web Speech API} dans Chrome / Edge pour un aperçu instantané ; (2)~\textit{MediaRecorder} puis transcription \textbf{Whisper} (OpenAI) lorsque le navigateur ne fournit pas Web Speech (Firefox inclus), avec boutons \textbf{Terminé} / \textbf{Annuler} et détection automatique FR / AR. L'import de fichier audio emprunte la même voie Whisper.

\section{Rôles et espace cabinet}

L'\textbf{admin cabinet} a le même atelier que l'avocat, plus l'équipe, le plan, le logo et le parrainage. L'\textbf{avocat} dicte, rédige, valide et exporte, sans payer. L'\textbf{assistant} ouvre le menu \textbf{Plaidoiries validées} : aperçu, consultation et téléchargement Word (.docx) / PDF des plaidoiries au statut \textit{validated} ou \textit{finalized} ; il dépose des pièces.

L'assistant n'utilise pas l'atelier de dictée, ni les fonctions d'IA de rédaction, de synthèse documentaire, ou de fiche orale d'audience (la \textbf{clé de barre} : document de préparation à plaider, généré depuis le dossier --- ce n'est ni un mot de passe, ni un code d'accès).

L'\textbf{admin plateforme} (WebVoice) supervise les cabinets depuis \texttt{/admin}, sans lire le contenu des affaires.

Le menu cabinet (dock fixe en bas sur mobile) : Plaidoirie, Dossiers, Documents, Planning, Activité, Équipe (admin cabinet). Paramètres : cabinet, plan et usage, sécurité, zone sensible. L'admin plateforme LexHub a un dock dédié (Aperçu, Cabinets, Plans, Leads, Logs, Coach) et un en-tête marque \textbf{LexHub} --- sans lien Paramètres cabinet.

\section{Expérience multi-support (desktop et mobile)}
\label{sec:mobile}

LexHub est un SaaS \textbf{web} unique (\url{https://lexhub.space}) : il n'existe pas d'application native App Store / Play Store séparée. L'atout est la \textbf{continuité d'usage} : le même compte, les mêmes rôles et le même contrôle d'accès sur \textbf{ordinateur} (rédaction approfondie) et sur \textbf{smartphone} (continuité hors cabinet).

\paragraph{Quel type de mobile ?}
Il ne s'agit pas d'une liste de marques, mais d'une adaptation aux \textbf{écrans tactiles courants} (largeurs typiques $360$--$430$\,px : iPhone, Galaxy, etc.), validée en navigateur mobile et en émulation (outils développeur). Au-delà d'environ $880$\,px de largeur, l'interface retrouve le mode bureau (dock latéral).

\paragraph{Impact UX.}
Sur téléphone : navigation en \textbf{dock fixé en bas} (gestes au pouce), mise en page mono-colonne, en-tête allégé. Les rôles avocat, assistant et admin plateforme conservent leurs spécificités (pas de dictée pour l'assistant ; pas de lecture des affaires pour l'admin plateforme).

\paragraph{Valeur métier.}
L'avocat peut reprendre un brouillon ou consulter un dossier en déplacement ; l'assistant récupère une plaidoirie validée ; l'exploitation LexHub reste pilotable depuis un mobile. Desktop et mobile ne sont pas deux produits : c'est \textbf{une seule expérience cabinet}, adaptée à la forme de l'écran.

\section{Inscription et entrée en équipe}

L'inscription d'un \textbf{nouveau cabinet} : e-mail, OTP à six chiffres, nom du cabinet, mot de passe. Le cabinet est créé \textbf{actif} (plan Découverte) ; la connexion est immédiate. Il n'y a pas de validation LexHub à l'inscription.

Pour rejoindre un cabinet \emph{existant}, l'admin envoie une \textbf{invitation par e-mail}. Le destinataire ouvre le lien, choisit son propre mot de passe, et accède au cabinet. Il n'existe plus de code cabinet circulant dans l'équipe.

\section{Les cinq étapes de la plaidoirie}
\label{sec:guide}

Libellés de production (rail « Étape $X$ / $5$ ») --- figures~\ref{fig:plaidoirie} et~\ref{fig:transcription}.

\begin{enumerate}[leftmargin=1.4em]
  \item \textbf{Dossier \& pièces.} Dossier obligatoire. Fiche d'identité juridique. À la création, chaque champ affiche une \textbf{aide contextuelle} (pourquoi le renseigner, exemple). Ajout de pièces. Ces données nourrissent la rédaction ; elles n'inventent pas le fond.
  \item \textbf{Dictée.} Reconnaissance vocale du navigateur (Web Speech, Chrome / Edge) ou capture micro \textbf{Whisper} (MediaRecorder --- Firefox inclus) ou import audio. Commandes démarrer / arrêter. La langue parlée (FR / AR) peut être détectée automatiquement.
  \item \textbf{Transcription.} Texte brut éditable. Bouton \textbf{Corriger} : forme standard, sans résumer ; les fautes de frappe / STT sont approximées intelligemment (pas de refus pour un texte encore compréhensible). Un texte sans phrase (touches au hasard) est \textbf{refusé} : l'IA n'invente pas une affaire. Surlignage orange des substitutions STT et alertes de contradiction. Avant la rédaction : feuille \textbf{Clarification rapide} (QCM 0--5) : LexHub anonymise dossier + dictée, le LLM prépare les questions, l'avocat confirme ou passe ; les réponses alimentent la composition. Le bouton \textbf{Nouvelle plaidoirie} propose de reprendre un brouillon local ou d'ouvrir une page vide (plus de réouverture silencieuse).
  \item \textbf{Rédaction IA.} Document éditable (structure \texttt{أولاً} \ldots \texttt{خامساً}, \texttt{لهذه الأسباب} / الملتمسات). Panneau \textbf{Assistant IA --- modifications} : historique type chat ; consignes ou questions en langage naturel (fautes OK) ; l'IA \textbf{répond pédagogiquement} (ex.\ rôle du titre « pour ces motifs ») \textbf{et} met à jour le corps (remplissage des sections vides à partir des faits déjà présents). UX micro alignée sur les assistants grand public : bandeau d'écoute, boutons \textbf{Terminé} / \textbf{Annuler}, détection auto FR/AR sans changer l'UI ; \textbf{Annuler la requête} / \textbf{Modifier la requête} pendant ou après l'envoi. L'avocat enregistre, valide ou assigne.
  \item \textbf{Export officiel Word (.docx) et PDF.} Génération du document Word officiel (standard modifiable des juridictions marocaines) et du PDF homologue : typographie Noto Naskh Arabic RTL, logo ou armoiries centrés, en-tête bilingue. Le .docx reste 100\,\% éditable avant audience. Après export réussi : statut \textit{finalized} (mise à jour UI). Les alertes « infos manquantes » de la rédaction ne s'affichent plus à cette étape.
\end{enumerate}

\section{Validation, assignation et assistant}

\begin{table}[H]
\centering
\caption{Collaboration autour d'une plaidoirie (comportement déployé)}
\label{tab:collab}
\small
\begin{tabular}{@{}lp{9.8cm}@{}}
\toprule
\textbf{Action} & \textbf{Effet réel} \\
\midrule
Valider & Statut \textit{validated} : \textbf{tous} les assistants du cabinet la voient. \\
Assigner & Valide le brouillon si besoin et la met en avant pour l'assistant choisi. \\
Exporter Word (.docx) / PDF & Statut \textit{finalized} : classée « Finalisée » (distinct du brouillon local). \\
Épingler & Mise en tête de liste. \\
Dévalider & Retour brouillon : plus visible des assistants. \\
\bottomrule
\end{tabular}
\end{table}

Trois statuts en production : \textbf{brouillon} (\texttt{draft}), \textbf{validée} (\texttt{validated}), \textbf{finalisée} (\texttt{finalized}). Valider pour l'assistant et finaliser après export sont deux gestes distincts.

L'assistant juridique est un collaborateur du cabinet, pas le moteur d'IA. Son droit de lecture des plaidoiries dépend du \textbf{statut validé ou finalisé} du document, non du fait d'être nommé dans une assignation. L'assignation oriente le travail ; elle ne constitue pas, à elle seule, le contrôle d'accès.

\section{Sécurité applicative}
\label{sec:securite}

La sécurité de LexHub combine des contrôles \emph{d'accès}, d'\emph{isolation} et d'\emph{exploitation}.

\paragraph{Accès.}
Chaque requête authentifiée porte un JWT. Les routes sont filtrées par rôle (avocat, assistant, admin cabinet, admin plateforme). L'inscription d'un cabinet passe par un OTP e-mail. Les invitations d'équipe sont des liens à usage nominatif, expirables, distincts du parrainage. Une authentification à deux facteurs (TOTP) est optionnelle.

\paragraph{Anti-robot (Cloudflare Turnstile).}
Turnstile est un widget Cloudflare qui distingue un humain d'un script, sans puzzle visuel du type CAPTCHA classique. Sur LexHub il n'est \textbf{pas} affiché à la première connexion : le formulaire e-mail / mot de passe suffit. Après \textbf{trois} échecs (seuil configuré), le frontend charge le widget ; l'API refuse le login tant que le jeton n'est pas validé auprès de Cloudflare. Le mécanisme est \textbf{actif en production} (clés réelles, pas les clés de test). Il ne s'applique pas à l'inscription (OTP e-mail) ni aux parcours métier (dictée, plaidoirie). En environnement \texttt{production}, le backend refuse de démarrer sans ces clés.

\paragraph{Confidentialité avant IA.}
Avant tout appel LLM (composition, clarification QCM, révision), les noms des parties sont remplacés par des jetons (\texttt{PARTIE\_A}, \ldots) puis restaurés dans le document présenté à l'avocat. L'audio Whisper et les pièces OCR transitent chiffrés (HTTPS) ; le stockage objet (MinIO) n'est pas exposé publiquement.

\paragraph{Isolation.}
Toutes les entités métier sont rattachées à un \texttt{cabinet\_id}. L'API refuse une ressource hors locataire. Le navigateur n'ouvre jamais PostgreSQL : il ne parle qu'à FastAPI en HTTPS.

\paragraph{Données au repos.}
Fichiers dans le magasin d'objets ; base relationnelle sur le VPS, ports liés à l'hôte local. Adminer, lorsqu'il est utilisé, n'écoute que sur \texttt{127.0.0.1} et s'atteint par tunnel SSH.

\paragraph{Abus.}
Limitation de débit sur l'authentification et les endpoints sensibles, complétée par Turnstile après échecs répétés. Journal d'activité. L'admin plateforme suspend ou purge un cabinet sans ouvrir les dossiers. Le journal d'audit plateforme enregistre l'adresse IP de la requête (traçabilité des connexions, pratique standard des SaaS) ; ce n'est pas une donnée métier. Les factures visibles côté plateforme sont les reçus de chaque cabinet (essai, packs), pas un livre de comptes interne WebVoice.

\paragraph{Paramètres cabinet.}
Outre l'abonnement, l'espace Paramètres porte la charte (logo, en-tête), la sécurité (2FA TOTP), les invitations d'équipe, le parrainage, le digest e-mail des notifications et l'export du planning au format \texttt{ICS} (Google Agenda / Outlook).

\paragraph{Parrainage.}
\label{sec:parrainage}
Le code de parrainage n'est pas un secret d'accès. Sa révélation ne compromet pas les dossiers ; elle permet au plus l'attribution d'un bonus d'essai, plafonnée.
